% !TeX program = pdflatex
\documentclass[11pt]{article}

\usepackage[margin=1in]{geometry}
\usepackage{lmodern}
\usepackage{microtype}
\usepackage{setspace}
\onehalfspacing

\usepackage{amsmath,amssymb,amsfonts,mathtools,bm,booktabs,tabularx,array}
\usepackage{enumitem}
\usepackage{hyperref}
\hypersetup{colorlinks=true,linkcolor=black,citecolor=black,urlcolor=black}

% ---- Unified notation ----
\newcommand{\R}{\mathbb{R}}
\newcommand{\E}{\mathbb{E}}
\newcommand{\Var}{\operatorname{Var}}
\newcommand{\Cov}{\operatorname{Cov}}
\newcommand{\tr}{\operatorname{tr}}
\newcommand{\diag}{\operatorname{diag}}
\newcommand{\vect}[1]{\bm{#1}}
\newcommand{\mat}[1]{\bm{#1}}
\newcommand{\zeros}[1]{\vect 0_{#1}}
\newcommand{\ones}[1]{\vect 1_{#1}}
\newcommand{\ind}{\mathbf{1}}
\newcommand{\Normal}{\mathcal{N}}
\newcommand{\IG}{\operatorname{IG}}
\newcommand{\StudentT}{\operatorname{t}}
\newcommand{\argmin}{\operatorname*{arg\,min}}

% ---- Supplement numbering ----
\newcounter{suppalgorithm}
\renewcommand{\thesuppalgorithm}{S\arabic{suppalgorithm}}
\newenvironment{suppalgorithm}[1]{%
  \refstepcounter{suppalgorithm}%
  \par\medskip\noindent\textbf{Algorithm~\thesuppalgorithm. #1.}\par\nobreak
}{\par\medskip}

\renewcommand{\theequation}{S\arabic{equation}}
\renewcommand{\thetable}{S\arabic{table}}
\renewcommand{\thesection}{S\arabic{section}}
\renewcommand{\thesubsection}{S\arabic{section}.\arabic{subsection}}
\renewcommand{\thesubsubsection}{S\arabic{section}.\arabic{subsection}.\arabic{subsubsection}}

\title{Supplementary Note: Exact Gaussian DESN Readouts under a Scaled Ridge Prior}
\author{Antonio Aguirre \and Raquel Prado \and Bruno Sans\'o}
\date{\today}

\begin{document}
\maketitle

\begin{abstract}
\noindent
This note derives the exact conjugate Gaussian, or normal-likelihood, deep echo
state network (DESN) readout used as a fast baseline and initialization device
for Q--DESN models. Conditional on the reservoir construction and washout, the
DESN states define a deterministic design matrix, so the Gaussian readout is a
Bayesian linear regression model on fixed nonlinear features. We use a scaled
ridge prior, \(\vect\beta\mid\omega^2\sim\Normal_d(\vect b,\omega^2\mat P^{-1})\),
and an inverse-gamma prior on the homoskedastic observation variance
\(\omega^2\). This model has an exact Normal--inverse-gamma posterior, a
closed-form marginal Student posterior for the readout coefficients, a
closed-form posterior predictive distribution for fixed future DESN features,
and an exact direct posterior sampler. Since the posterior is available in
closed form, no Gibbs sampler, CAVI/VB algorithm, or ELBO approximation is
needed for this Gaussian baseline. Regularized-horseshoe variants are mentioned
only as conditionally conjugate alternatives, not as part of the exact baseline.
\end{abstract}

\section{Scope and Relation to Q--DESN}
\label{sec:scope}

After reservoir construction, spectral scaling, layer reduction, and washout,
the DESN feature map is fixed. Posterior inference for the readout is therefore
conditional on a deterministic design matrix
\(\mat X=[\vect x_1^\top;\ldots;\vect x_T^\top]\). The Gaussian DESN considered
in this note replaces the AL or exAL working likelihood with the homoskedastic
Normal likelihood
\begin{equation}
\label{eq:gaussian_likelihood_scope}
\vect y\mid \vect\beta,\omega^2,\mat X
\sim \Normal_T(\mat X\vect\beta,\omega^2\mat I_T),
\qquad \omega^2>0.
\end{equation}
The variance parameter is denoted by \(\omega^2\), not \(\sigma\), to avoid
confusion with the AL/exAL scale parameter in Q--DESN.

The Gaussian DESN is not a quantile model. If the Gaussian likelihood is taken
literally, \(\vect x_t^\top\vect\beta\) targets a conditional mean under
squared-error loss. In Q--DESN, by contrast, \(\vect x_t^\top\vect\beta\) is the
conditional \(p_0\)-quantile under an AL or exAL working likelihood. The role of
this note is therefore computational:
\begin{enumerate}[label=(\roman*)]
\item provide an exact mean-regression baseline on the same fixed DESN features;
\item provide a conditioning and regularization diagnostic for the readout;
\item provide stable initial values for AL and exAL Q--DESN algorithms;
\item avoid unnecessary approximate inference for a model whose posterior is
available analytically.
\end{enumerate}

\section{Fixed DESN Design and Notation}
\label{sec:fixed_design}

Let
\begin{equation}
\label{eq:y_X_beta_dimensions}
\vect y=(y_1,\ldots,y_T)^\top\in\R^T,
\qquad
\mat X=[\vect x_1^\top;\ldots;\vect x_T^\top]\in\R^{T\times d},
\qquad d=r+1.
\end{equation}
Each row vector \(\vect x_t=(1,\tilde{\vect x}_t^\top)^\top\in\R^d\) contains
an intercept and \(r\) fixed DESN readout features. The coefficient vector is
partitioned as
\begin{equation}
\label{eq:beta_partition}
\vect\beta=(\beta_0,\beta_1,\ldots,\beta_r)^\top\in\R^d,
\qquad
\vect\beta_{1:r}=(\beta_1,\ldots,\beta_r)^\top.
\end{equation}
The first coefficient is the intercept. The remaining \(r\) coefficients are
slopes associated with the fixed DESN features. No reservoir weight, input
weight, reducer, leak rate, or spectral-radius parameter is inferred in this
Gaussian readout.

For any \(\vect\beta\), define
\begin{equation}
\label{eq:residuals}
\vect e(\vect\beta)=\vect y-\mat X\vect\beta,
\qquad
S_y(\vect\beta)=\|\vect y-\mat X\vect\beta\|^2.
\end{equation}
The identity matrix of size \(m\) is \(\mat I_m\).

\section{Distributional Conventions}
\label{sec:dist_conventions}

All inverse-gamma distributions use the reciprocal-scale rate parameterization
\begin{equation}
\label{eq:ig_density}
p(x\mid a,b)=\frac{b^a}{\Gamma(a)}x^{-a-1}\exp(-b/x),
\qquad x>0,
\end{equation}
written \(x\sim\IG(a,b)\). If \(x\sim\IG(a,b)\), then
\begin{equation}
\label{eq:ig_moments}
\E(x^{-1})=\frac{a}{b},
\qquad
\E(x)=\frac{b}{a-1}\quad(a>1),
\qquad
\E(\log x)=\log b-\psi(a),
\end{equation}
where \(\psi(\cdot)\) is the digamma function. For a Gaussian factor,
\begin{equation}
\label{eq:normal_entropy}
H\{\Normal_d(\vect m,\mat S)\}
=\frac12\{d(1+\log 2\pi)+\log|\mat S|\}.
\end{equation}
The notation \(\StudentT_\nu(\vect m,\mat S)\) denotes a multivariate Student
\(t\) distribution with degrees of freedom \(\nu\), location \(\vect m\), and
scale matrix \(\mat S\). Its covariance is \(\nu(\nu-2)^{-1}\mat S\), for
\(\nu>2\).

\section{Exact Scaled-Ridge Gaussian DESN}
\label{sec:scaled_ridge}

\subsection{Model}
\label{subsec:model}

The exact Gaussian baseline uses the conjugate scaled ridge prior
\begin{align}
\vect y\mid\vect\beta,\omega^2,\mat X
&\sim \Normal_T(\mat X\vect\beta,\omega^2\mat I_T),
\label{eq:likelihood}\\
\vect\beta\mid\omega^2
&\sim \Normal_d(\vect b,\omega^2\mat P^{-1}),
\label{eq:scaled_prior}\\
\omega^2&\sim\IG(a_\omega,b_\omega),
\label{eq:omega_prior}
\end{align}
where \(a_\omega,b_\omega>0\), \(\mat P\in\R^{d\times d}\) is positive
definite, and
\begin{equation}
\label{eq:b_P_definition}
\vect b=(b_0,0,\ldots,0)^\top,
\qquad
\mat P=\diag(P_0,\kappa_\beta^{-2},\ldots,\kappa_\beta^{-2}).
\end{equation}
If
\(\beta_0\mid\omega^2\sim\Normal(b_0,\omega^2V_0)\), then
\(P_0=V_0^{-1}\). More general positive definite \(\mat P\) can be used without
changing the derivation.

The scaling by \(\omega^2\) in \eqref{eq:scaled_prior} is the key conjugacy
condition. It is not identical to an independent ridge prior
\(\vect\beta\sim\Normal_d(\vect b,\mat P^{-1})\). The scaled version is retained
here because it yields an exact one-shot posterior and is used only as a
baseline and initialization device for Q--DESN.

\subsection{Full joint distribution}
\label{subsec:joint}

The joint density factors as
\begin{equation}
\label{eq:joint_factorization}
p(\vect y,\vect\beta,\omega^2\mid\mat X)
=p(\vect y\mid\vect\beta,\omega^2,\mat X)
 p(\vect\beta\mid\omega^2)p(\omega^2).
\end{equation}
Writing all terms that depend on unknowns,
\begin{align}
\label{eq:joint_density_explicit}
p(\vect y,\vect\beta,\omega^2\mid\mat X)
&=(2\pi\omega^2)^{-T/2}
\exp\left\{-\frac{1}{2\omega^2}
(\vect y-\mat X\vect\beta)^\top(\vect y-\mat X\vect\beta)\right\}
\nonumber\\
&\quad\times
(2\pi\omega^2)^{-d/2}|\mat P|^{1/2}
\exp\left\{-\frac{1}{2\omega^2}
(\vect\beta-\vect b)^\top\mat P(\vect\beta-\vect b)\right\}
\nonumber\\
&\quad\times
\frac{b_\omega^{a_\omega}}{\Gamma(a_\omega)}
(\omega^2)^{-a_\omega-1}\exp(-b_\omega/\omega^2).
\end{align}
Equivalently, up to constants independent of \((\vect\beta,\omega^2)\),
\begin{align}
\label{eq:log_joint}
\log p(\vect y,\vect\beta,\omega^2\mid\mat X)
&=\frac12\log|\mat P|+a_\omega\log b_\omega-\log\Gamma(a_\omega)
\nonumber\\
&\quad
-\left(a_\omega+1+\frac{T+d}{2}\right)\log\omega^2
\nonumber\\
&\quad
-\frac{1}{2\omega^2}\left[
\|\vect y-\mat X\vect\beta\|^2
+(\vect\beta-\vect b)^\top\mat P(\vect\beta-\vect b)
+2b_\omega\right].
\end{align}

\subsection{Completing the square}
\label{subsec:complete_square}

Define
\begin{align}
\mat P_n&=\mat P+\mat X^\top\mat X,
\label{eq:Pn}\\
\vect h_n&=\mat P\vect b+\mat X^\top\vect y,
\label{eq:hn}\\
\vect m_n&=\mat P_n^{-1}\vect h_n,
\label{eq:mn}\\
a_n&=a_\omega+\frac{T}{2},
\label{eq:an}\\
B_n&=b_\omega+\frac12\left(
\vect y^\top\vect y+
\vect b^\top\mat P\vect b-
\vect m_n^\top\mat P_n\vect m_n
\right).
\label{eq:Bn}
\end{align}
Then
\begin{align}
\label{eq:complete_square_identity}
&\|\vect y-\mat X\vect\beta\|^2
+(\vect\beta-\vect b)^\top\mat P(\vect\beta-\vect b)
\nonumber\\
&\qquad=
(\vect\beta-\vect m_n)^\top\mat P_n(\vect\beta-\vect m_n)
+\vect y^\top\vect y+
\vect b^\top\mat P\vect b-
\vect m_n^\top\mat P_n\vect m_n.
\end{align}
Substituting \eqref{eq:complete_square_identity} into the joint density gives
an exact Normal--inverse-gamma posterior.

\subsection{Exact posterior}
\label{subsec:exact_posterior}

The posterior factorization is
\begin{align}
\omega^2\mid\vect y,\mat X&\sim\IG(a_n,B_n),
\label{eq:omega_posterior}\\
\vect\beta\mid\omega^2,\vect y,\mat X
&\sim\Normal_d(\vect m_n,\omega^2\mat P_n^{-1}).
\label{eq:beta_cond_posterior}
\end{align}
Equivalently,
\begin{equation}
\label{eq:joint_posterior_kernel}
p(\vect\beta,\omega^2\mid\vect y,\mat X)
=p(\vect\beta\mid\omega^2,\vect y,\mat X)p(\omega^2\mid\vect y,\mat X).
\end{equation}
The marginal coefficient posterior is multivariate Student \(t\):
\begin{equation}
\label{eq:beta_student_posterior}
\vect\beta\mid\vect y,\mat X
\sim
\StudentT_{2a_n}\left(\vect m_n,\frac{B_n}{a_n}\mat P_n^{-1}\right).
\end{equation}
Hence
\begin{align}
\E(\vect\beta\mid\vect y,\mat X)&=\vect m_n,
\label{eq:posterior_mean_beta}\\
\Var(\vect\beta\mid\vect y,\mat X)
&=\frac{B_n}{a_n-1}\mat P_n^{-1},
\qquad a_n>1,
\label{eq:posterior_var_beta}\\
\E(\omega^2\mid\vect y,\mat X)&=\frac{B_n}{a_n-1},
\qquad a_n>1.
\label{eq:posterior_mean_omega}
\end{align}
The posterior mode of \(\vect\beta\) and its posterior mean coincide at
\(\vect m_n\), while the marginal posterior for \(\vect\beta\) carries the
additional uncertainty induced by integrating out \(\omega^2\).

\section{Posterior Predictive Distributions}
\label{sec:predictive}

For a new fixed DESN feature vector \(\vect x_*\in\R^d\), the posterior
predictive distribution is
\begin{equation}
\label{eq:predictive_univariate}
y_*\mid\vect y,\mat X,\vect x_*
\sim
\StudentT_{2a_n}\left(
\vect x_*^\top\vect m_n,
\frac{B_n}{a_n}\{1+\vect x_*^\top\mat P_n^{-1}\vect x_*\}
\right).
\end{equation}
More generally, for a fixed future design matrix \(\mat X_*\in\R^{H\times d}\),
\begin{equation}
\label{eq:predictive_multivariate}
\vect y_*\mid\vect y,\mat X,\mat X_*
\sim
\StudentT_{2a_n}\left(
\mat X_*\vect m_n,
\frac{B_n}{a_n}\{\mat I_H+\mat X_*\mat P_n^{-1}\mat X_*^\top\}
\right).
\end{equation}
The fixed-design condition matters. In multi-step DESN forecasting, future
features may depend on previously simulated responses. Conditional on any
particular generated feature path, \eqref{eq:predictive_multivariate} applies.
If the feature path is recursively simulated, posterior prediction should
propagate draws through the reservoir recursion rather than treating
\(\mat X_*\) as known in advance.

\section{Exact Log Marginal Likelihood}
\label{sec:log_marginal}

Integrating out both \(\vect\beta\) and \(\omega^2\) gives the closed-form
marginal likelihood
\begin{align}
\label{eq:log_marginal}
\log p(\vect y\mid\mat X)
&=-\frac{T}{2}\log(2\pi)
+\frac12\log|\mat P|-\frac12\log|\mat P_n|
+a_\omega\log b_\omega-a_n\log B_n
\nonumber\\
&\quad+\log\Gamma(a_n)-\log\Gamma(a_\omega).
\end{align}
This replaces any need for an ELBO in the Gaussian baseline. It is also a
strong implementation check: exact posterior samplers and numerical posterior
summaries should be consistent with \eqref{eq:log_marginal} and
\eqref{eq:omega_posterior}--\eqref{eq:beta_student_posterior}.

\section{Exact Direct Posterior Sampler}
\label{sec:direct_sampler}

Because the posterior is available exactly, no Markov chain is required.
Independent posterior draws are generated directly as follows.

\begin{suppalgorithm}{Direct sampler for the exact scaled-ridge Gaussian DESN}
\label{alg:direct_sampler}
\emph{Input:} fixed design \(\mat X\), response \(\vect y\), prior
\((\vect b,\mat P,a_\omega,b_\omega)\), and desired number of draws \(M\).
\begin{enumerate}[label=(\arabic*)]
\item Compute \(\mat P_n\), \(\vect m_n\), \(a_n\), and \(B_n\) from
\eqref{eq:Pn}--\eqref{eq:Bn}.
\item For \(m=1,\ldots,M\):
  \begin{enumerate}[label=(\alph*)]
  \item draw \(\omega^{2(m)}\sim\IG(a_n,B_n)\);
  \item draw
  \(\vect\beta^{(m)}\mid\omega^{2(m)}
  \sim\Normal_d(\vect m_n,\omega^{2(m)}\mat P_n^{-1})\).
  \end{enumerate}
\item For fixed future \(\mat X_*\), optionally draw
\(\vect y_*^{(m)}\sim\Normal_H(\mat X_*\vect\beta^{(m)},
\omega^{2(m)}\mat I_H)\), or use the analytic Student distribution in
\eqref{eq:predictive_multivariate}.
\end{enumerate}
\end{suppalgorithm}

Algorithm~\ref{alg:direct_sampler} is a direct posterior sampler, not a Gibbs
sampler. A Gibbs sampler would only reproduce a posterior that is already
available in closed form. Similarly, CAVI/VB would only approximate an exact
posterior. For this reason, Gibbs, CAVI, and ELBO derivations are intentionally
omitted from the main Gaussian baseline.

\section{Use as Initialization for AL and exAL Q--DESN}
\label{sec:initialization}

The posterior mean \(\vect m_n\) gives a stable Gaussian readout estimate on the
same fixed DESN features used by Q--DESN. For a target quantile level
\(p_0\in(0,1)\), use it as a mean-regression warm start and then shift the
intercept to target the desired empirical quantile.

Define Gaussian residuals
\begin{equation}
\label{eq:init_residuals}
r_t^N=y_t-\vect x_t^\top\vect m_n,
\qquad t=1,\ldots,T.
\end{equation}
Let
\begin{equation}
\label{eq:empirical_resid_quantile}
c_{p_0}=\widehat Q_{p_0}(r_1^N,\ldots,r_T^N)
\end{equation}
be the empirical \(p_0\)-quantile of the Gaussian residuals. Set
\begin{align}
\beta_0^{\mathrm{init}}(p_0)&=m_{n,0}+c_{p_0},
\label{eq:init_beta0}\\
\beta_j^{\mathrm{init}}(p_0)&=m_{n,j},
\qquad j=1,\ldots,r.
\label{eq:init_betaj}
\end{align}
Then the empirical \(p_0\)-quantile of
\(y_t-\vect x_t^\top\vect\beta^{\mathrm{init}}(p_0)\) is approximately zero.
This makes \(\vect x_t^\top\vect\beta^{\mathrm{init}}(p_0)\) a more appropriate
initial quantile path than the raw Gaussian mean path.

With the check loss
\begin{equation}
\label{eq:check_loss}
\rho_p(u)=u\{p-\ind(u<0)\},
\end{equation}
initialize the AL/exAL scale by
\begin{equation}
\label{eq:init_sigma}
\sigma^{\mathrm{init}}
=\frac1T\sum_{t=1}^T
\rho_{p_0}\left(y_t-\vect x_t^\top\vect\beta^{\mathrm{init}}(p_0)\right)
+\epsilon,
\end{equation}
where \(\epsilon>0\) is a small numerical floor. For AL, no shape parameter or
\(s_t\) variables are used. For exAL, initialize at the AL special case,
\begin{equation}
\label{eq:init_gamma}
\gamma^{\mathrm{init}}=0,
\end{equation}
or at a small jittered value near zero if the implementation uses an open
transformed support. Stable latent starts are
\begin{equation}
\label{eq:init_latents}
v_t^{\mathrm{init}}=\sigma^{\mathrm{init}},
\qquad
s_t^{\mathrm{init}}=\sqrt{2/\pi}\quad\text{for exAL only}.
\end{equation}

For variational Q--DESN implementations, a Gaussian covariance scale can also
be initialized from
\begin{equation}
\label{eq:init_covariance}
\widehat{\Var}(\vect\beta\mid\vect y,\mat X)
=\frac{B_n}{a_n-1}\mat P_n^{-1},
\qquad a_n>1,
\end{equation}
or from the conditional covariance
\(\widehat\omega^2\mat P_n^{-1}\) with
\(\widehat\omega^2=B_n/(a_n-1)\). These covariance quantities should be viewed
as numerical initializers rather than posterior uncertainty statements for the
AL/exAL quantile model.

\section{Omitted Variants and Why They Are Not Main Derivations}
\label{sec:omitted_variants}

\subsection{Independent, unscaled ridge prior}
\label{subsec:unscaled_omitted}

The independent ridge prior
\begin{equation}
\label{eq:unscaled_prior_omitted}
\vect\beta\sim\Normal_d(\vect b,\mat P^{-1}),
\qquad
\omega^2\sim\IG(a_\omega,b_\omega)
\end{equation}
is conditionally conjugate: \(\vect\beta\mid\omega^2,\vect y,\mat X\) is
Gaussian and \(\omega^2\mid\vect\beta,\vect y,\mat X\) is inverse-gamma.
However, it does not yield the same one-shot Normal--inverse-gamma posterior as
Section~\ref{sec:scaled_ridge}, because the coefficient prior is not scaled by
\(\omega^2\). It would require a Gibbs sampler or a mean-field approximation
for posterior computation. Since the purpose of this note is an exact Gaussian
baseline and initializer, the unscaled prior is not developed further here.

\subsection{Regularized horseshoe readout prior}
\label{subsec:rhs_omitted}

A Gaussian DESN with the Nishimura--Suchard regularized-horseshoe
\texttt{rhs\_ns} prior is also conditionally conjugate under its product
representation. Given the local, global, and slab scales, the coefficient block
is Gaussian, and the scale blocks can have inverse-gamma updates. Nevertheless,
after integrating over those shrinkage scales, the marginal posterior is not a
single closed-form Normal--inverse-gamma or Student distribution. Thus the
\texttt{rhs\_ns} Gaussian model is useful as a conditional Gibbs or CAVI warm
start, but it is not the exact closed-form baseline derived in this note.

For Q--DESN initialization, the recommended default is therefore
\begin{equation}
\label{eq:init_pipeline}
\begin{gathered}
\text{exact scaled-ridge Gaussian DESN}\\
\longrightarrow\ \text{quantile-shifted }\vect\beta^{\mathrm{init}}(p_0)\\
\longrightarrow\ \text{AL/exAL Q--DESN}.
\end{gathered}
\end{equation}
If adaptive RHS shrinkage is required in the final Q--DESN fit, RHS scale
initializations can still be constructed from the shifted ridge coefficients,
for example
\begin{align}
\tau^{\mathrm{init}}
&=\operatorname{median}_{j=1,\ldots,r}|\beta_j^{\mathrm{init}}|+
\epsilon,
\label{eq:init_tau_rhs}\\
\lambda_j^{\mathrm{init}}
&=\max\left\{\lambda_{\min},
\frac{|\beta_j^{\mathrm{init}}|}{\tau^{\mathrm{init}}}\right\},
\qquad j=1,\ldots,r,
\label{eq:init_lambda_rhs}\\
\zeta^{\mathrm{init}}
&=\max_{j=1,\ldots,r}|\beta_j^{\mathrm{init}}|+
\epsilon.
\label{eq:init_zeta_rhs}
\end{align}
These are initialization heuristics, not an exact RHS posterior calculation.

\section{Implementation Notes}
\label{sec:implementation}

The main computational object is the positive definite precision matrix
\(\mat P_n=\mat P+\mat X^\top\mat X\). Implementations should use Cholesky
factorization of \(\mat P_n\) for all solves, log determinants, and Gaussian
sampling. Avoid explicit matrix inversion except in small diagnostic scripts.
For example,
\begin{equation}
\label{eq:solve_not_inverse}
\vect m_n=\mat P_n^{-1}\vect h_n
\end{equation}
should be computed by solving \(\mat P_n\vect m_n=\vect h_n\).

For the log marginal likelihood,
\(\log|\mat P_n|\) and \(\log|\mat P|\) should be obtained from Cholesky factors.
For posterior covariance diagnostics, use triangular solves to compute selected
diagonal entries or linear-function variances rather than forming a dense
inverse when \(d\) is large.

Feature preprocessing should be fixed before inference. Centering and scaling
non-intercept DESN features can improve conditioning, but the intercept column
should remain an explicit column of ones. The prior precision \(P_0\) controls
intercept regularization separately from the slope precision
\(\kappa_\beta^{-2}\).

\section{Validation and Consistency Checks}
\label{sec:validation}

\subsection{Dimension checks}

\begin{enumerate}[label=(\alph*)]
\item \(\mat X\in\R^{T\times d}\), \(\mat X^\top\mat X\in\R^{d\times d}\),
\(\mat P\in\R^{d\times d}\), and \(\mat P_n\in\R^{d\times d}\).
\item \(\vect h_n=\mat P\vect b+\mat X^\top\vect y\in\R^d\), hence
\(\vect m_n\in\R^d\).
\item \(B_n\) is scalar and positive when the prior is proper and
\(\mat P_n\) is positive definite.
\item \(\vect x_*^\top\mat P_n^{-1}\vect x_*\) is scalar and nonnegative.
\end{enumerate}

\subsection{Conjugacy checks}

\begin{enumerate}[label=(\alph*)]
\item The scaled prior contributes \(-d\log\omega^2/2\) to the joint density,
but this factor is canceled by the Gaussian integral over \(\vect\beta\) when
obtaining \(\omega^2\mid\vect y\). Therefore the marginal shape is
\(a_\omega+T/2\), not \(a_\omega+(T+d)/2\).
\item Conditional on \(\omega^2\), the posterior precision for \(\vect\beta\) is
\(\omega^{-2}\mat P_n\), and the posterior covariance is
\(\omega^2\mat P_n^{-1}\).
\item Integrating \(\omega^2\) gives the Student posterior in
\eqref{eq:beta_student_posterior} with degrees of freedom \(2a_n\).
\item The posterior predictive variance includes both observation noise and
coefficient uncertainty through
\(1+\vect x_*^\top\mat P_n^{-1}\vect x_*\).
\end{enumerate}

\subsection{Numerical checks}

\begin{enumerate}[label=(\alph*)]
\item Independent draws from Algorithm~\ref{alg:direct_sampler} should match
\(\E(\omega^2\mid\vect y,\mat X)=B_n/(a_n-1)\) and
\(\E(\vect\beta\mid\vect y,\mat X)=\vect m_n\), within Monte Carlo error.
\item The empirical covariance of sampled \(\vect\beta\) should match
\(B_n(a_n-1)^{-1}\mat P_n^{-1}\), for \(a_n>1\).
\item The direct Student sampler for \(\vect\beta\) and the two-step
\(\omega^2,\vect\beta\mid\omega^2\) sampler should give the same marginal
summaries.
\item The log determinant terms in \eqref{eq:log_marginal} should be computed
from Cholesky factors and checked for finite values.
\item The condition numbers of \(\mat X^\top\mat X\), \(\mat P\), and
\(\mat P_n\) should be monitored. Poor conditioning indicates a feature-scaling
or reservoir-design issue, not a posterior-derivation issue.
\end{enumerate}

\subsection{Q--DESN initialization checks}

\begin{enumerate}[label=(\alph*)]
\item After the intercept shift in \eqref{eq:init_beta0}, the empirical
\(p_0\)-quantile of
\(y_t-\vect x_t^\top\vect\beta^{\mathrm{init}}(p_0)\) should be close to zero.
\item The initialized AL/exAL scale in \eqref{eq:init_sigma} must be strictly
positive.
\item The initialized exAL shape value must lie in the quantile-dependent
support \((L,U)\). If \(0\) is represented only as a boundary case in code, use
a small interior jitter around zero.
\item If RHS shrinkage scales are initialized by
\eqref{eq:init_tau_rhs}--\eqref{eq:init_zeta_rhs}, all scales must be finite and
positive.
\end{enumerate}

\subsection{Final self-audit}

This note keeps only the scaled-ridge Gaussian DESN as the main model because
it is the version with an exact closed-form posterior. It does not present
Gibbs, CAVI, or ELBO algorithms for the main model because those methods are
unnecessary approximations or redundant samplers for an analytically available
posterior. Unscaled ridge and \texttt{rhs\_ns} priors are described only to
clarify why they are not the exact baseline: both are useful conditionally
conjugate models, but neither yields the same one-shot closed-form marginal
posterior after all nuisance quantities are integrated out.

\end{document}
